\scanpage{410}
\mspage{324}
% archive sheet number: 409

% 1–5
\deleted{par} $\overline{\unsure{j|J}}_\alpha$, $u_\alpha$ par le composé $\overline{u}_\alpha:\varphi(i_{\unsure{\delta|\alpha}})\longrightarrow j_\alpha\to\overline{j}_\alpha$, et $\tau'_{\alpha,\alpha_0}$ par le composé $\overline{\tau'}_{\alpha,\alpha_0}:j_{\alpha_0}\to j_\alpha\to\overline{j}_{\unsure{\alpha_0|\alpha}}$, (sans toucher à $i_\alpha,v_{\alpha_0},\tau_{\alpha_0,\alpha}$), on peut obtenir la commutativité voulue, d'où \deleted{$\unsure{3_\alpha}$} $3_\alpha(\beta)$.

% 6–7
Reste la contrainte $4_\alpha(\beta)$. Pour ceci, considérons \unsure{le|les} \unsure{diagramme|diagrammes} \unsure{composite|composites}
% 8–10
\diagram{\begin{tikzcd}[ampersand replacement=\&,column sep=large]
j_\beta\arrow[r,"v_\beta"]\arrow[d,"{\tau'_{\alpha_0,\beta}}"'] \& \varphi(i_{\beta+1})\arrow[d,"{\varphi(\tau_{\alpha,\beta+1})}"]\\
j_{\alpha_0}\arrow[r,"v_{\alpha_0}"] \& \varphi(i_{\alpha_0\deleted{\unsure{\alpha}}+1})=\varphi(i_\alpha)\quad .
\end{tikzcd}}

% 11–17
\deleted{\unsure{Pour tout}} Si $\alpha_0$ n'est pas un ordinal limite, \deleted{mais} de la forme $\alpha_1+1$, on peut encore se ramener au cas où $\beta=\alpha_1$, en procédant comme ci-dessus. Mais on \deleted{peut} \unsure{procédera} dans tous les cas, \deleted{par une} sans un tel artifice (!), en considérant les doubles-flèches
% 18
\diagram{\begin{tikzcd}[ampersand replacement=\&,column sep=huge]
(\unsure{*})\qquad j_\beta\arrow[r,shift left=2.5pt,"{\varphi(\tau_{\alpha,\beta+1})v_\beta}"]\arrow[r,shift right=2.5pt,"{v_{\alpha_0}\tau'_{\alpha_0,\beta}}"'] \& \varphi(i_\alpha)
\end{tikzcd}}
% 19–20
déduites du carré précédent. Comme \deleted{$J'$ \unsure{est} $J$ est \unsure{grand} devant $\unsure{\pi|\Pi}$.}
